\documentclass[12pt]{article}
\usepackage[utf8]{inputenc}
\usepackage[T1]{fontenc}
\usepackage[margin=1in]{geometry}
\usepackage{setspace}
\usepackage{xcolor}
\usepackage[colorlinks=true, urlcolor=blue!60!black, linkcolor=blue!60!black]{hyperref}
\setlength{\parindent}{0em}
\setlength{\parskip}{0.6em}
\renewcommand{\baselinestretch}{1.0}
\begin{document}
\begin{center}
{\LARGE \textit{Jones v. Alfred H. Mayer Co.}, 392 U.S. 409 (1968)\par}
\vspace{0.5em}
{\large Case Brief --- Two-Page Edition\par}
\end{center}

\section{Memory Jogger}
Section 1982 of the 1866 Civil Rights Act bars private racial discrimination in property sales, a valid use of Congress's Thirteenth Amendment power over the ``badges and incidents of slavery'' (\href{https://supreme.justia.com/cases/federal/us/392/409/}{392 U.S.\ 409}).

\section{Detailed Case Facts}
In 1965 a private developer refused to sell Joseph Lee Jones, a Black man, and his white wife a home in its St.\ Louis County subdivision solely because Mr.\ Jones was Black. They sued under 42 U.S.C.\ \S\ 1982, which guarantees all citizens ``the same right'' to buy property ``as is enjoyed by white citizens.'' The Fair Housing Act of 1968 passed while the case was pending but did not cover their claim.

\section{Procedural History}
The district court dismissed, holding that \S\ 1982 reaches only state action; the Eighth Circuit (Blackmun, J.) affirmed.

\section{Judicial Votes}
Reversed, 7--2. Stewart wrote for Warren, Black, Douglas, Brennan, Fortas, and Marshall; Douglas also concurred; Harlan, joined by White, dissented.

\section{Holding}
Section 1982 bars all racial discrimination, private as well as public, in property transactions. The Thirteenth Amendment has no state-action limit, and Section 2 lets Congress rationally identify and eliminate the badges and incidents of slavery. \textit{Hodges v.\ United States} was overruled to the extent inconsistent.

\section{Analysis of Opinions}
Stewart held that the Fair Housing Act left \S\ 1982 intact and that the 1866 Congress meant to protect freedmen against private ``custom, or prejudice'' as well as the Black Codes. Under the \textit{Civil Rights Cases}, the Thirteenth Amendment permits legislation reaching private individuals, and Congress could rationally find that exclusion from white communities, which ``herds men into ghettos,'' was a relic of slavery. Douglas catalogued slavery's persisting badges. Harlan, dissenting, read ``the same right'' as equal legal capacity against state laws, not a right to compel private sellers, and would have dismissed the writ given the new Fair Housing Act.

\section{Examples: Future Applications}
\textbf{Same-Side:} a private school's refusal to admit Black students violates \S\ 1981 (\textit{Runyon v.\ McCrary}). \textbf{Opposite-Side:} a landlord's religious or sex discrimination falls outside \S\ 1982. \textbf{Fence-Sitter:} a Section 2 statute banning sexual-orientation or disability discrimination in housing.

\section{Critique}
Critics such as Gerhard Casper and Senator Ervin fault the history, and scholars warn of tension with \textit{City of Boerne}. Defenders such as Darrell Miller and William Carter show the 1866 Congress targeted private ``white cartels'' and praise the revival of the Thirteenth Amendment.

\section{Key Quotations}
\begin{quote}
``Surely Congress has the power under the Thirteenth Amendment rationally to determine what are the badges and the incidents of slavery, and the authority to translate that determination into effective legislation.''
\end{quote}
\noindent --- Justice Stewart, 392 U.S.\ at 440.

\end{document}
